\documentclass[../../../include/open-logic-section]{subfiles}

\begin{document}

\olfileid{sth}{card-arithmetic}{ch}

\olsection{د پيوستار فرضيه}

تېرې پايلې (په سمه توګه) دې ته اشاره وکړه چې د اصلي عددونو توان
به ډېر \emph{اسان} و، که نامتناهي اصلي عددونه د $2$ له توانونو سره
«په نېغه توګه» چلند وکړي؛ يعنې، د \olref[opps]{lem:SizePowerset2Exp}
له مخې، د توان سټونو له اخيستلو سره۔ خو دا نه شو فرضولے چې
نامتناهي اصلي عددونه په رښتيا له توان سټونو سره داسې نېغ چلند کوي۔

د دې خبرې د سپړلو د پيل دپاره څو ګټورې نښې راوړو۔

\begin{defn}
چې $ \cardsucc{\cardfont{a}} $ له $\cardfont{a}$ څخه په کلکه لوے
تر ټولو کوچنے اصلي عدد دے، لاندې دوه نامتناهي لړۍ تعريفوو:
\begin{align*}
	\aleph_{0} &\defis \omega & 		
	\beth_{0} &\defis \omega\\
	\aleph_{\alpha \ordplus 1} &\defis \cardsucc{(\aleph_{\alpha})} &
	\beth_{\alpha+1} &\defis \cardexpo{2}{\beth_{\alpha}}\\
	\aleph_{\alpha} &\defis \bigcup_{\beta< \alpha} \aleph_{\beta} &
	\beth_{\alpha} &\defis \bigcup_{\beta < \alpha}\beth_{\beta} & \text{چې $\alpha$ حدي ترتيبي عدد وي}.
	\end{align*}
\end{defn}

د $\cardsucc{\cardfont{a}}$ تعريف پر ځاے دے، ځکه
\olref[cardinals][classing]{lem:NoLargestCardinal} وايي چې د هر اصلي
عدد $\cardfont{a}$ دپاره له $\cardfont{a}$ څخه لوے يو اصلي عدد
شته، او د نامتناهيت هاخوا
استقرا دا تضمينوي چې له $\cardfont{a}$ څخه لوے يو \emph{تر ټولو کوچنے}
اصلي عدد شته۔ د دغو لړيو پاتې تعريفونه د نامتناهيت هاخوا بازګښت
له لارې ورکول کېږي۔
\textbf{د سرچينې يادښت (OLSTH-020):} چاپي جمله دلته د
$\cardfont{a}$ د پاتې تعريف خبره کوي؛ په سياق کښې د دغو دوو لړيو
پاتې تعريفونه مراد دي۔

کانتور د «$\aleph$» نښه راوړه؛ دې ته \emph{اَلېف} وايي۔ دا د
عبراني الفبا لومړے تورى او د «نامتناهي» د عبراني کلمې لومړے
تورے دے۔ پيرس د «$\beth$» نښه راوړه؛ دې ته \emph{بېت} وايي،
چې د عبراني الفبا دويم تورى دے۔\footnote{پيرس دا نښه د ۱۹۰۰م کال
د دسمبر په يوه ليک کښې کانتور ته کارولې وه۔ خو په هماغه ليک
کښې ئې د دې دپاره ناسم استدلال هم وکړ چې $\beth_\alpha$ د
$\alpha \geq \omega$ دپاره شتون نه لري۔}
اوس دا نښې موږ ته نامتناهي اصلي عددونه راکوي۔

\begin{prop}
$\aleph_\alpha$ او $\beth_\alpha$ د هر ترتيبي عدد $\alpha$
دپاره اصلي عددونه دي۔
\end{prop}

\begin{proof}
دواړه پايلې د نامتناهيت هاخوا په ساده استقرا ثابتېږي۔
$\aleph_0 = \beth_0 = \omega$ د
\olref[cardinals][classing]{omegaisacardinal} له مخې اصلي عدد دے۔
که فرض کړو چې $\aleph_\alpha$ او $\beth_\alpha$ دواړه اصلي
عددونه دي، نو $\aleph_{\alpha+1}$ او
$\beth_{\alpha+1}$ په څرګند ډول د اصلي عددونو په توګه تعريف
شوي دي۔ او د اصلي عددونو د يوه سټ اتحاد د
\olref[cardinals][classing]{unioncardinalscardinal} له مخې اصلي
عدد دے۔
\end{proof}
\noindent
سربېره پر دې، هر نامتناهي اصلي عدد يو $\aleph$ دے:

\begin{prop}
که $\cardfont{a}$ نامتناهي اصلي عدد وي، نو
$\cardfont{a} = \aleph_\gamma$، د يوه يوازيني $\gamma$ دپاره۔
\end{prop}

\begin{proof}
پر اصلي عددونو د نامتناهيت هاخوا استقرا کوو۔ د استقرا دپاره فرض
کړئ چې که $\cardfont{b} < \cardfont{a}$ وي، نو
$\cardfont{b} = \aleph_{\gamma_\cardfont{b}}$۔ که
$\cardfont{a} = \cardsucc{\cardfont{b}}$ د کوم
$\cardfont{b}$ دپاره وي، نو $\cardfont{a} =
\cardsucc{(\aleph_{\gamma_\cardfont{b}})}=
\aleph_{\gamma_\cardfont{b}+1}$۔ که $\cardfont{a}$ د هېڅ
اصلي عدد تالي نه وي، نو ځکه چې اصلي عددونه ترتيبي عددونه دي،
$\cardfont{a} = \bigcup_{\cardfont{b} < \cardfont{a}} \cardfont{b} =
\bigcup_{\cardfont{b} < \cardfont{a}}{\aleph_{\gamma_\cardfont{b}}}$؛
نو $\cardfont{a} = \aleph_\gamma$، چېرته چې $\gamma =
\bigcup_{\cardfont{b} < \cardfont{a}}\gamma_\cardfont{b}$۔
\end{proof}
\textbf{د سرچينې يادښت (OLSTH-021):} چاپي ثبوت د کوچنيو اصلي
عددونو په فرض کښې متناهي عددونه هم رانغاړي، حال دا چې د هغو
دپاره د اَلېف شاخص نه دے تعريف شوے؛ د لومړي نامتناهي اصلي عدد
بنسټيز حالت هم بېل نه دے ثابت شوے۔ پورته فورمولونه
د سرچينې په بڼه پاتې دي او د دې ثبوت تشه لا پرانيستې ده۔

څرنګه چې هر نامتناهي اصلي عدد يو $\aleph$ دے، دا پوښتنه راولاړېږي:
ايا هر نامتناهي اصلي عدد يو $\beth$ هم دے؟ که داسې \emph{وے}،
نو نامتناهي اصلي عددونو به د توان سټونو د اخيستلو له عمل سره
«په نېغه توګه» چلند کاوه۔ په رښتيا، بيا به مو لاندې ادعا درلوده:

\begin{defish}
\emph{د پيوستار عمومي فرضيه} (GCH)۔
$\aleph_\alpha  = \beth_\alpha$، د هر $\alpha$ دپاره۔
\end{defish}

سربېره پر دې، که GCH سمه وے، نو د اصلي عددونو توان مو د پام وړ
ساده کولے شو۔ په ځانګړې توګه، ثابتولے مو شول چې کله
$\cardfont{b} < \cardfont{a}$ وي، د
$\cardexpo{\cardfont{a}}{\cardfont{b}}$ اندازه د
$\cardfont{a}\leq \cardexpo{\cardfont{a}}{\cardfont{b}} \leq
\cardsucc{\cardfont{a}}$ تر منځ راځي۔ بيا مو دقيق شرطونه هم
ورکولے شول چې له دغو دوو امکانونو کوم يو پېښېږي، يعنې
$\cardfont{a} = \cardexpo{\cardfont{a}}{\cardfont{b}}$
او که $\cardexpo{\cardfont{a}}{\cardfont{b}} =
\cardsucc{\cardfont{a}}$۔\footnote{دا شرط د
\emph{هم‌پايۍ} له مخې ټاکل کېږي۔}

خو GCH يوه \emph{فرضيه} ده، \emph{قضيه} نه ده۔ په رښتيا،
\citet{Godel1938} ثابته کړه چې که $\ZFC$ سازګاره وي، نو
$\ZFC + \text{GCH}$ هم سازګاره ده۔ خو وروسته څرګنده شوه
چې $\lnot$GCH هم له $\ZFC$ سره زياتولے شو۔ د GCH تر ټولو
ساده غير بديهي \emph{ځانګړے حالت} دا دے:

\begin{defish}
\emph{د پيوستار فرضيه} (CH)۔ $\aleph_1 = \beth_1$۔
\end{defish}

\citet{Cohen1963} ثابته کړه چې که $\ZFC$ سازګاره وي، نو
$\ZFC + \lnot\text{CH}$ هم سازګاره ده۔ نو د پيوستار فرضيه
له $\ZFC$ څخه خپلواکه ده۔

دې ته د پيوستار فرضيه ځکه ويل کېږي چې «پيوستار» د حقيقي عددونو
د کرښې، $\Real$، بل نوم دے۔
\olref[opps]{continuumis2aleph0} وايي چې $\card{\Real} =
\beth_1$۔ نو د پيوستار فرضيه وايي چې د طبيعي عددونو د سټ د
اصلي شمېر، $\aleph_0 = \beth_0$، او د پيوستار د اصلي شمېر،
$\beth_1$، تر منځ بل اصلي عدد نشته۔

له $\ZFC$ څخه د (G)CH د \emph{خپلواکۍ} له مخې، د دې د
\emph{رښتياوالي} په باب څه ووايو؟ په دې باب \emph{ډېرې}
خبرې شته۔ د سټونو په نظريه کښې د وروستيو څېړنو ډېره ګټوره
برخه هم د دغو پوښتنو څېړلو ته متوجه ده۔ خو دوه لنډ ټکي
خامخا د يادولو وړ دي۔

لومړے: له دغو رسمي خپلواکۍ پايلو څخه دا \emph{سمدستي}
نه راځي چې د GCH يا CH د رښتيا ارزښت \emph{ناتاکلے}
دے۔ کېدای شي يوازې نورو داسې اصلونو ته اړتيا ولرو چې موږ
ته طبيعي ښکاري او پوښتنه په يوه يا بله خوا هواره کړي۔
ګېډل پخپله همدا غبرګون مناسب ګاڼه۔

دويم: له $\ZFC$ څخه د CH خپلواکي بېشکه \emph{حيرانوونکې}
ده، خو په لفظي معنا \emph{ناباورېدونکې} نه ده۔ خبره دا ده چې
$\ZFC$، تر کومه ځايه چې موږ ته وايي، د يوه اصلي عدد څخه د
هغه تالي ته تګ ښايي د توان سټونو له اخيستلو څخه ډېر نازک
عمل وغواړي۔
 %The operation of taking powersets moves us from one stage of the
 %hierarchy to its successor stage (from $V_{\alpha}$ to
 %$V_{\alpha+1}$). 

له دې دوو خبرو وروسته، که نور معلومات غواړئ، اوس به د دې
بحث د بېلابېلو دريځونو لرونکو فيلسوفانو او رياضي‌پوهانو
ليکنې ولولئ۔\footnote{کېدای شي لوستل د \citet[\S15.6]{Potter2004}
څخه پيل کړئ۔}

\end{document}
